\scanpage{9}
\uline{Thm 2} Pour $\forall\ \pi$ discret, il existe un $\pi$ fibré universel,
i.e. un espace classifiant (défini à isom unique près dans la cat. homotopique ..)

Soit $X$ un espace sur lequel $\pi$ opère à \unsure{g.}
et prenons
% Second pass: after "opère" the ink shows a dotted stroke and an S-shaped sign with a short dash; the "⊃ X"
% of the next line stands just below them. Hua batch-01 (p. 8) leaves the word illegible and reads the group
% of signs as "⊃ X^δ ?". Not settled.

\diagram{\begin{tikzcd}[ampersand replacement=\&,column sep=large]
\&E_\pi\times_\pi X\overset{\deleted{\unsure{déf|df|E}}}{=}\underbrace{(E_\pi\times X)/\pi}\arrow[dl]\arrow[dr]\&\unsure{\supset}X\\
B_\pi=E_\pi/\pi\&\&X/\underset{\deleted{E_\pi}}{\pi}
\end{tikzcd}}
% Second pass: faint pencil marks at the left edge, level with the space between "et prenons" and the lower
% sequence (an arrow-like stroke and a short stroke); no letters can be read.

\diagram{\begin{tikzcd}[ampersand replacement=\&,column sep=small]
\pi_2(B_\pi)\arrow[r]\arrow[d,equal]\&\pi_1(X)\arrow[r,dash]\&\pi_1(X,\pi)\arrow[r]\&\pi\arrow[r]\&\pi_0(\unsure{V|X})\arrow[d,equal]\\
0\&\&\&\&1
\end{tikzcd}}

\mspage{8}
